\newpage
\section{Архитектура и реализация системы}
\label{ch:impl}

В этой главе описано, как обзорные решения главы~1 превращаются в работающую систему. Изложение
идёт по ходу данных: от сбора корпуса и построения эмбеддингов через кластеризацию и метакластеры
к рекомендательному пайплайну и веб-приложению. Сравнительные эксперименты, обосновывающие выбор
конкретных алгоритмов и значений гиперпараметров, вынесены в главу~3 — здесь приводятся уже
выбранные варианты.

\subsection{Общая архитектура}
\label{sec:arch}

Система построена как последовательность независимых этапов обработки данных и логически делится
на две фазы (рис.~\ref{fig:arch}). \textbf{Предподсчёт и обучение} выполняются заранее, офлайн:
корпус OpenAlex кодируется эмбеддингами, понижается по размерности, кластеризуется, темы получают
интерпретируемое представление и объединяются в метакластеры; параллельно по истории публикаций
обучается модель представления автора, и для всех авторов заранее вычисляются эмбеддинги.
\textbf{Инференс} происходит онлайн, в ответ на запрос: по выбранному автору быстро извлекаются
кандидаты, языковая модель их оценивает, отбирает, ищет доказательства и формирует объяснение.
Разделение существенно по нагрузке — тяжёлые вычисления (эмбеддинги, кластеризация, обучение)
делаются один раз, а на запрос пользователя приходится лишь лёгкий retrieval и несколько вызовов
языковой модели.

\begin{figure}[ht]
\centering
\resizebox{\textwidth}{!}{%
\begin{tikzpicture}[
  >={Stealth[round]},
  font=\scriptsize,
  box/.style={rectangle, rounded corners, draw=black!55, align=center,
              minimum height=8mm, inner sep=2.5pt, fill=blue!5},
  data/.style={box, fill=black!7},
  arr/.style={->, draw=black!65, thick}
]
% --- Предподсчёт ---
\node[data] (n1) at (0,0)      {OpenAlex\\(AI-корпус)};
\node[box]  (n2) at (3.1,0)    {Эмбеддинги\\e5-base};
\node[box]  (n3) at (6.3,0)    {UMAP 10D\\+ HDBSCAN};
\node[box]  (n4) at (9.6,0)    {c-TF-IDF, MMR,\\LLM-метки};
\node[box]  (n5) at (12.8,0)   {Метакластеры\\(HDBSCAN)};
\node[data] (n6) at (3.1,-1.8) {История\\авторов};
\node[box]  (n7) at (6.3,-1.8) {Transformer /\\GraphSAGE};
\node[box]  (n8) at (9.6,-1.8) {Эмбеддинги\\авторов};
\draw[arr] (n1)--(n2); \draw[arr] (n2)--(n3); \draw[arr] (n3)--(n4); \draw[arr] (n4)--(n5);
\draw[arr] (n6)--(n7); \draw[arr] (n7)--(n8);
\draw[arr] (n2)--(n6);
\begin{scope}[on background layer]
\node[draw=black!40, dashed, rounded corners, fill=black!3, inner sep=8pt,
      fit=(n1)(n5)(n6)(n8),
      label={[anchor=north west, font=\scriptsize\itshape]north west:Предподсчёт и обучение}] (offline) {};
\end{scope}
% --- Инференс ---
\node[data] (m1) at (0,-4.7)    {Запрос\\автора};
\node[box]  (m2) at (3.1,-4.7)  {Dense\\retrieval};
\node[box]  (m3) at (6.3,-4.7)  {LLM-оценка\\и отбор};
\node[box]  (m4) at (9.6,-4.7)  {Доказательства\\и объяснение};
\node[box]  (m5) at (12.8,-4.7) {Веб-\\интерфейс};
\draw[arr] (m1)--(m2); \draw[arr] (m2)--(m3); \draw[arr] (m3)--(m4); \draw[arr] (m4)--(m5);
\begin{scope}[on background layer]
\node[draw=black!40, dashed, rounded corners, fill=green!4, inner sep=8pt,
      fit=(m1)(m5),
      label={[anchor=north west, font=\scriptsize\itshape]north west:Инференс}] (online) {};
\end{scope}
\draw[arr, dashed] (offline.south) -- node[right, font=\scriptsize]{артефакты предподсчёта} (online.north);
\end{tikzpicture}%
}
\caption{Архитектура системы: фаза предподсчёта и обучения (офлайн) и фаза инференса (онлайн)}
\label{fig:arch}
\end{figure}

Переиспользуемый код вынесен в пакет \code{src/}, разбитый по доменам (сбор данных, эмбеддинги,
кластеризация, метакластеры, рекомендации, оценка, веб); запускаемые сценарии — тонкие обёртки в
\code{scripts/run_*.py}, а гиперпараметры зафиксированы в \code{configs/*.yaml}.

\paragraph{Воспроизводимость.} Результаты системы детерминированы. Эмбеддинги статей вычисляются
один раз и далее переиспользуются; все источники случайности зафиксированы единым сидом
($\mathrm{seed} = 42$ для UMAP, K-Means и обучения модели, тогда как HDBSCAN детерминирован при
фиксированном входе); температура языковой модели выставлена в $0$. При неизменных входных данных
и коде повторный запуск даёт тот же результат.

\paragraph{Технологический стек.} Выбор инструментов подчинён задаче, и каждый пункт обоснован:

\begin{itemize}
  \item \textbf{Эмбеддинги статей} — multilingual-e5-base: многоязычный корпус и быстрый локальный
        инференс на 509\,447 статей;
  \item \textbf{Понижение размерности} — UMAP: сохраняет локальную структуру многообразия лучше
        линейных методов;
  \item \textbf{Кластеризация} — HDBSCAN: не требует заранее задавать число тем и умеет выделять
        шум;
  \item \textbf{Представление автора} — Transformer-энкодер и GraphSAGE (PyTorch): агрегируют
        историю публикаций и граф прошлых соавторств;
  \item \textbf{Языковая модель} — OpenAI gpt-4.1-mini: разметка тем, оценка релевантности и
        генерация объяснений;
  \item \textbf{Backend} — FastAPI; \textbf{frontend} — Deck.gl: отрисовка сотен тысяч точек
        карты в браузере.
\end{itemize}

\subsection{Построение эмбеддингов статей}
\label{sec:emb}

Каждая статья превращается в один вектор, отражающий её смысл. На вход берётся заголовок и
аннотация; к ним добавляется префикс \texttt{query:}, требуемый моделью
multilingual-e5-base~\cite{wang2024multilinguale5}. Единый префикс для всех статей гарантирует,
что эмбеддинги получены в одном режиме и сравнимы между собой.

Текст токенизируется с ограничением $\mathrm{max\_length} = 512$ токенов и кодируется моделью
\code{intfloat/multilingual-e5-base}; вектор L2-нормализуется — тогда его скалярное произведение
с другим вектором равно косинусной близости.

Корпус — 509\,447 статей OpenAlex по подполю Artificial Intelligence; полный набор фильтров
отбора описан в разделе~3.1. Результат этапа — матрица эмбеддингов \code{embeddings.npy}
размера $(509\,447, 768)$ и список идентификаторов \code{paper_ids.json}. Список задаёт канонический
порядок статей: по нему все последующие этапы сопоставляют строку матрицы с конкретной статьёй.

\subsection{Кластеризация научных публикаций}
\label{sec:cluster}

Тематическая структура строится конвейером в стиле BERTopic: эмбеддинги $\to$ UMAP $\to$ HDBSCAN
$\to$ c-TF-IDF $\to$ MMR $\to$ разметка языковой моделью. UMAP понижает размерность с 768 до
10 измерений, сохраняя локальную структуру, после чего HDBSCAN выделяет плотные
тематические сгустки, не требуя заранее задавать число тем и помечая разреженные статьи как шум.

Каждая тема описывается набором характерных слов. Сначала по кластерам считается c-TF-IDF (кластер
рассматривается как один большой документ), затем MMR отбирает из топ-слов разнообразный, не
дублирующий сам себя набор терминов. Наконец, языковая модель получает эти термины и заголовки
представительных статей и возвращает короткое человекочитаемое название темы. Примеры получившихся
названий fine-тем: \emph{Memory-Enabled AI Chatbots for Personalized Support},
\emph{Physics-Informed Neural Networks for PDE Inference},
\emph{Deep Learning for Solar Irradiance Forecasting},
\emph{Human-Guided Reinforcement Learning in Robotics}.

Качество вариантов оценивается набором \emph{независимых} метрик (без сводного взвешенного балла):

\begin{itemize}
  \item \emph{когерентность} $c$ — насколько статьи внутри темы относятся к одному
        исследовательскому направлению; оценивается языковой моделью по шкале 0–3 и взвешивается
        по размеру темы; в сравнительных таблицах нормируется к $[0,1]$ делением на 3;
  \item \emph{различимость} $d$ — доля пар тем, которые языковая модель сочла дублирующими; чем
        меньше, тем лучше разделены темы;
  \item \emph{сбалансированность размеров} — отсутствие «мегатем», поглощающих корпус;
        измеряется по доле статей, приходящейся на пять крупнейших кластеров (top5, меньше — лучше);
  \item \emph{доля шума} $n$ — доля статей, не отнесённых ни к одной плотной теме (свойственна
        только HDBSCAN);
  \item \emph{энтропия назначений} $e$ — для мягких методов (приписывают статью нескольким темам
        с вероятностями) высокая энтропия означает размытые границы тем.
\end{itemize}

По совокупности этих метрик выбран вариант \texttt{hdbscan\_fine} с \textbf{979} темами
($\mathrm{min\_cluster\_size} = 10$, $\mathrm{min\_samples} = 5$, метод отбора кластеров — leaf):
он даёт наибольшую когерентность ($c = 2{,}80$, нормированная $\hat{c} = 0{,}93$) при приемлемой
различимости ($d = 0{,}11$). Цена выбора — высокая доля шума ($n = 0{,}65$): значительная часть
статей не попадает в плотные кластеры, что ожидаемо для неоднородного AI-корпуса; тем не менее
детализация и интерпретируемость fine-слоя важнее для карты. Сравнительные таблицы, бейзлайны
и обоснование выбора приведены в разделе~3.2.

\subsection{Метакластеры}
\label{sec:meta}

979 fine-тем — слишком детально для обзора предметной области целиком, поэтому темы дополнительно
объединяются в \textbf{76} метакластеров. Это \emph{семантическое} объединение, а не геометрическая
иерархия: каждая тема представляется метадокументом (её LLM-название, топ-термины c-TF-IDF и
заголовки представительных статей), метадокумент кодируется тем же эмбеддером
multilingual-e5-base, и уже эти эмбеддинги тем кластеризуются HDBSCAN.

Когерентность и различимость метакластеров оценены полноценной LLM-оценкой с отдельными
промптами для мета-уровня (судья оценивает, образуют ли fine-темы одну \emph{широкую}
исследовательскую область, а не узкую fine-тему). По когерентности лучшим оказался вариант
\texttt{meta\_hdbscan\_medium} ($K = 76$): $\hat{c} = 0{,}90$ ($c = 2{,}71$), доля дублирующих пар
$0\,\%$. На этом уровне HDBSCAN помечает около $11\,\%$ тем как разреженные; такие темы
присоединяются к ближайшему по вероятности метакластеру, так что все 979 тем в итоге распределены
по 76 группам. Полная сравнительная таблица — в разделе~3.2.

При $K = 76$ каждая группа объединяет в среднем $\approx 13$ fine-тем; подпись метакластера
наследует LLM-название доминирующей темы в группе и поэтому звучит как конкретное поднаправление
AI, а не как абстрактная метка верхнего уровня. Примеры названий метакластеров:
\emph{Evolutionary Algorithms and Optimization Methods},
\emph{Quantum Computing with Superconducting Circuits},
\emph{Deep Learning for Medical Image Analysis},
\emph{Interpretable and Explainable NLP Models}.

Метакластеры — это крупный смысловой слой над темами, но \textbf{не} области на карте, и причина
техническая. Метакластеры получены кластеризацией кластеров, то есть статьи и темы живут в разных
векторных пространствах. Двумерная UMAP-проекция способна хорошо разделить либо статьи, либо темы,
но не то и другое сразу; наша карта построена так, чтобы разделять \emph{статьи}, и области
метакластеров на ней неизбежно накладываются друг на друга. Поэтому метакластеры используются как
подписи и контекст, а от их геометрической отрисовки на карте сознательно отказались.

\subsection{Рекомендательная система: представление автора и извлечение кандидатов}
\label{sec:retrieval}

Рекомендатель состоит из обучаемого этапа извлечения (этот раздел) и этапа переранжирования
языковой моделью (раздел~\ref{sec:rerank}). Обучается только первый.

\paragraph{Представления автора.} В работе сравниваются три способа получить вектор автора.
Первый — \code{mean_author_embedding}: L2-нормализованное среднее эмбеддингов последних статей
до cutoff. Это простой и сильный semantic baseline без обучения. Второй — Transformer-retriever
(\code{recommendation/model.py}): он получает последовательность эмбеддингов статей автора,
упорядоченных по времени; в cluster-aware вариантах к каждой статье дополнительно присоединяется
soft membership $q$ по темам. Базовый вариант без кластеров использует $q$ нулевой размерности,
fine-вариант использует $q = 980$ (979 fine-кластеров + шум), meta-вариант агрегирует эти признаки
до 76 метакластеров + шум.

В Transformer-retriever каждая статья проецируется в скрытое пространство
$d_{\text{model}} = 256$. Чтобы модель учитывала темп публикаций, разрыв во времени между соседними
статьями (в годах) разбивается на интервалы 0, 1, 2, 3 и 4+ лет, и каждому интервалу сопоставлен
обучаемый вектор-добавка к представлению статьи. К последовательности добавляется специальный
обучаемый токен автора (по аналогии с \texttt{[CLS]}), чьё выходное представление
$h_{\mathrm{author}}$ агрегирует историю.

Предсказанный вектор интересов автора $\hat{e}$ живёт в том же пространстве, что и эмбеддинги
статей, потому что retrieval сравнивает авторов по косинусной близости их векторов. Модель строит
$\hat{e}$ как \emph{остаток} поверх среднего эмбеддинга истории:
\begin{equation}
\label{eq:author-emb}
\hat{e} = \mathrm{normalize}\!\left(\bar{h} + \alpha \cdot \mathrm{MLP}(h_{\mathrm{author}})\right),
\end{equation}
где $\bar{h}$ — среднее эмбеддингов статей автора, $\mathrm{MLP}$ — двухслойная голова с GELU,
$\alpha$ — обучаемый масштаб остатка, инициализированный нулём. При $\alpha = 0$ модель совпадает
с mean-бейзлайном (среднее эмбеддингов истории); обучение добавляет лишь полезное отклонение от
него. Такой якорь важен: mean-бейзлайн напрямую использует качественные paper embeddings, и без
остаточной схемы обучаемый retriever проигрывал ему.

\paragraph{Графовый вариант.} Третий способ — GraphSAGE-retriever
(\code{recommendation/graph_retrieval.py}). Узлы графа — 35\,954 квалифицированных автора, ребро
между авторами проводится, если они имели совместную публикацию до cutoff 2024. В графе
464\,752 направленных author-author ребра; веса построчно нормируются, повторные старые
коллаборации не увеличивают вес ребра. Начальный признак узла $x_v$ — L2-нормализованное среднее
эмбеддингов последних статей автора, то есть тот же сигнал, что в mean-бейзлайне.

Один слой GraphSAGE вычисляет сообщения отдельно по типам рёбер:
\begin{equation}
\label{eq:graphsage-layer}
h_v^{(\ell+1)} =
\operatorname{LayerNorm}\!\left(
W_0 h_v^{(\ell)} +
\sum_{r \in \mathcal{R}} W_r
\sum_{u \in \mathcal{N}_r(v)} w_{vu}^{(r)} h_u^{(\ell)}
\right),
\end{equation}
где $h_v^{(0)} = x_v$, $\mathcal{R}$ — типы отношений, $\mathcal{N}_r(v)$ — соседи узла $v$ по
отношению $r$, $w_{vu}^{(r)}$ — нормированный вес ребра. В реализованной модели два слоя
($L = 2$), скрытая размерность 256, dropout 0.2. Выход снова строится как остаток поверх исходного
семантического признака:
\begin{equation}
\label{eq:graphsage-out}
z_v = \operatorname{normalize}\!\left(x_v + \beta W_{\mathrm{out}} h_v^{(L)}\right).
\end{equation}

Вариант с метакластерами не разрезает граф на независимые подграфы. Это один гетерогенный граф:
к авторским узлам добавлены 76 узлов метакластеров и отношения author-meta / meta-author. Связь
автора с метакластером имеет вес по суммарной массе его статей в метакластере; сохраняются веса
$\geq 0{,}02$, а при отсутствии таких весов берётся самый вероятный метакластер. В этом варианте
добавляется 172\,376 направленных author-meta/meta-author рёбер. Он проверяет, помогают ли
метакластеры как смысловые hub-узлы в графовом retrieval.

\paragraph{Обучение.} Модель обучается multi-positive \textbf{coauthor InfoNCE} на косинусной
близости. Для примера $i$ с предсказанным вектором $\hat{e}_i$ и множеством позитивов
$\mathcal{P}_i$ (будущие квалифицированные соавторы, не встречавшиеся в истории до cutoff) потеря
равна
\begin{equation}
\label{eq:loss}
\mathcal{L}_i = -\log \frac{\sum_{j \in \mathcal{P}_i} \exp\!\bigl(\cos(\hat{e}_i,\, e_j)/\tau\bigr)}
{\sum_{k \in \mathcal{N}_i} \exp\!\bigl(\cos(\hat{e}_i,\, e_k)/\tau\bigr)},
\end{equation}
где $e_j$ — эмбеддинг соавтора в том же варианте модели, $\mathcal{N}_i$ — сэмплированные
негативы из пула авторов, позитивы исключены из знаменателя (не конкурируют как negatives),
$\tau = 0{,}05$ — температура. Обе обучаемые семьи используют AdamW и early stopping по
val nDCG@50. Для Transformer задано $\mathrm{lr} = 2\cdot10^{-4}$ и
$\mathrm{weight\_decay} = 0{,}05$; для GraphSAGE — $\mathrm{lr} = 10^{-3}$ и
$\mathrm{weight\_decay} = 10^{-2}$. У Transformer скорость обучения линейно растёт первые 10\,\%
шагов, затем плавно снижается по косинусу.

Качество на обучении дополнительно контролируется метрикой nDCG@10 на валидации, где правильным
ответом считается фактически наблюдавшееся последующее соавторство (рис.~\ref{fig:ndcg}). Абсолютные
значения малы — это прямое следствие разреженности целевого сигнала, которая разбирается в
разделе~3.3.

\begin{figure}[ht]
\centering
\includegraphics[width=0.78\textwidth]{val_ndcg10.png}
\caption{Валидационная метрика nDCG@10 Transformer-представления автора по эпохам обучения}
\label{fig:ndcg}
\end{figure}

\paragraph{Извлечение кандидатов (dense retrieval).} После обучения для каждого
квалифицированного автора (всего их \textbf{35\,954} — тех, у кого $\geq 10$ статей) предсказывается
вектор интересов. Эти векторы L2-нормализуются; для запрошенного автора вычисляется косинусная
близость со всеми остальными, исключаются он сам и его существующие соавторы (рекомендуем
\emph{новые} связи), и берутся топ-$N$ ближайших (в основной схеме $N = 10$). Это и есть кандидаты для
переранжирования. Подчеркну: на этом этапе отбираются именно ближайшие по вектору исследователи —
с похожей тематической траекторией; содержательная оценка релевантности выполняется уже на
следующем этапе.

\subsection{Переранжирование и генерация объяснений на основе языковой модели}
\label{sec:rerank}

Dense retrieval даёт быстрый, но грубый список — он не учитывает содержательные детали. Эту работу
выполняет пятиступенчатый пайплайн на языковой модели \texttt{gpt-4.1-mini} (рис.~\ref{fig:arch},
фаза инференса).

\begin{enumerate}
  \item \textbf{Retrieval.} Топ-$N$ кандидатов из раздела~\ref{sec:retrieval}.
  \item \textbf{Оценка релевантности.} Каждый кандидат оценивается отдельным вызовом языковой
        модели по шкале 0–3 (см. ниже): в промпт подаются недавние статьи автора и кандидата.
        Вызовы независимы и выполняются параллельно.
  \item \textbf{Отбор.} Берутся топ-$K$ кандидатов (по умолчанию $K = 3$) по оценке.
  \item \textbf{Извлечение доказательств.} Для каждого отобранного кандидата модель находит не
        более двух сильнейших попарных связей «статья автора\,$\leftrightarrow$\,статья кандидата»,
        возвращая индексы связываемых статей (только из переданных в промпт), краткую конкретную
        причину и ссылки OpenAlex.
  \item \textbf{Объяснение.} Модель формирует связный текст рекомендаций, ссылаясь только на
        реально переданные статьи — выдумывать их запрещено.
\end{enumerate}

Шкала релевантности: \textbf{0} — нерелевантно; \textbf{1} — слабое/общее совпадение (только широкая
близость в духе «оба про AI»); \textbf{2} — релевантно (конкретная общая задача, метод, тип данных
или взаимодополняющая экспертиза); \textbf{3} — очень релевантно (сильное конкретное совпадение).
Каждому кандидату присваивается \emph{тип связи}: прямое тематическое совпадение, совпадение по
методу, взаимодополняющая экспертиза или слабое/общее совпадение.

Ключевое требование к качеству — опираться на \emph{конкретные} связи, а не на общие фразы. В
промптах прямо запрещены обоснования вида «both papers discuss AI methodologies»; действует правило
\emph{«если связь сводится лишь к широкой близости по AI/ML — не включать её»}. Кандидаты с типом
«слабое/общее совпадение» отфильтровываются из финального списка. Именно здесь, в отличие от
retrieval, система способна различать взаимодополняющую экспертизу — это делает суждение языковой
модели, а не близость векторов.

\subsection{Веб-приложение}
\label{sec:web}

Все этапы сведены в веб-приложение: backend на FastAPI раздаёт статику и JSON-API, frontend на
ванильном JavaScript и Deck.gl отрисовывает сотни тысяч точек. Основные маршруты API перечислены
в таблице~\ref{tab:api}.

\begin{table}[ht]
\caption{Основные маршруты API backend}
\label{tab:api}
\centering
\footnotesize
\begin{tabular}{ll}
\toprule
Маршрут & Назначение \\
\midrule
\texttt{GET /api/map/state} & Точки статей, центры и подписи тем и метакластеров \\
\texttt{GET /api/papers/\{id\}} & Детали статьи \\
\texttt{GET /api/clusters/info/\{id\}} & Сводка по теме или метакластеру \\
\texttt{GET /api/authors/search} & Автодополнение поиска авторов \\
\texttt{POST /api/recommendations/\{id\}} & Запуск пайплайна рекомендаций \\
\bottomrule
\end{tabular}
\end{table}

Главный экран — карта научного пространства, где каждая точка это статья, окрашенная по своей теме
(рис.~\ref{fig:map}). Навигация построена на зуме: на дальнем масштабе видны крупные подписи
метакластеров, на среднем проявляются названия тем, на глубоком — заголовки отдельных статей.
Подписи плавно проявляются и затухают по уровню зума, а их пересечения устраняются встроенной
коллизионной фильтрацией Deck.gl.

\begin{figure}[ht]
\centering
\includegraphics[width=\textwidth]{ui_map.png}
\caption{Карта научного пространства: каждая точка — статья, цвет наследуется от метакластера}
\label{fig:map}
\end{figure}

Клик по статье открывает боковую панель с деталями — заголовком, аннотацией, годом, авторами,
темой, метакластером и ссылкой на OpenAlex; поле поиска авторов поддерживает автодополнение с
транслитерацией кириллицы (рис.~\ref{fig:panel}). Выбор автора запускает пайплайн рекомендаций, и в
модальном окне показываются кандидаты с типом связи, кратким объяснением и компактными
доказательствами — попарными связями статей с короткой причиной и ссылками OpenAlex
(рис.~\ref{fig:recs}).

\begin{figure}[ht]
\centering
\includegraphics[width=0.95\textwidth]{ui_paper_panel.png}
\caption{Боковая панель статьи и поиск автора с автодополнением}
\label{fig:panel}
\end{figure}

\begin{figure}[ht]
\centering

\includegraphics[width=0.92\textwidth]{ui_recommendations.png}
\caption{Окно рекомендаций: кандидат, тип связи, объяснение и попарные доказательства}
\label{fig:recs}
\end{figure}

\FloatBarrier
\subsection{Выводы и результаты по главе}
\label{sec:ch2-summary}

В главе описана реализация всех подсистем. Корпус из 509\,447 статей Artificial Intelligence
представлен L2-нормализованными эмбеддингами multilingual-e5-base; конвейер в стиле BERTopic дал
979 интерпретируемых тем, объединённых в 76 метакластеров. Рекомендатель реализован в две стадии:
обучаемые Transformer- и GraphSAGE-представления автора (остаток поверх mean-бейзлайна,
coauthor InfoNCE) извлекают ближайших кандидатов, а пятиступенчатый пайплайн на языковой модели
переранжирует их и формирует объяснения с конкретными попарными доказательствами. Всё интегрировано в
веб-приложение с интерактивной картой и поиском. Количественные обоснования сделанных выборов —
сравнение кластеризации с бейзлайнами, retrieval-метрики и LLM-оценка качества — приводятся в главе~3.
